% generated by roam.py handbook export -- не редактировать, править заметки
% 20261004181419-метод_золотого_сечения.org (body 0d7f3fee7371)

\begin{method}[{Метод золотого сечения}]\label[method]{note:96DED7D9-3918-45C5-9850-E05030941645}
  Пусть $f$~--- \hyperref[note:CD52D818-AA25-471C-97C2-5DBE9787CE53]{унимодальна} на $[a;b]$;
  НУО $l[a;b]= 1$, где $l$~--- функция длины отрезка;

  Мы хотим выбрать такие две точки $x < y$ из этого отрезка, чтобы сократить отрезок, где лежит минимум, используя свойство: \hyperref[note:1D154EFA-62D8-4BC3-B9CD-CCDE31F166B0]{Монотонность унимодальной функции}.
  Далее будем  называть такие точки $x,y$~--- \emph{тестовыми}, так как в них мы и будем вычислять значение функции.

  Но мы хотим сделать что-то немного более хитрое.
  \textbf{Мы хотим считать функцию только 1 раз за итерацию, а не 2.}

  Для этого мы зададим ограничения на тестовые точки.

  \begin{enumerate}
    \item $l[a;x] = l[y;b]$ (Каждый раз сокращаем отрезок на одну величину).
      Обозначим $l[a;y] = K$, $l[a;x]= 1-K$. Это гарантирует скорость сходимости (мы всегда знаем размер оставшегося отрезка), иначе бы отрезки были бы неравными, и так как мы не контролируем, в какой отрезок мы пойдем, скорость бы определялась большим отрезком.
    \item Одна из тестовых точек текущего отрезка должна стать тестовой точкой нового отрезка (это и обеспечивает вычисление функции 1 раз). Пусть после обрезания остался отрезок $[a;y]$.  Тогда новый $y_1$ должен стоять от начала на расстоянии $K \cdot K$ (в следующей итерации, мы считаем, что новый отрезок длинны 1 и домножаем на $K$, чтобы найти расстояние до $y_1$), и мы хотим, чтобы он совпал с $x$.
  \end{enumerate}

  Исходя из этих ограничений мы имеем
  $1 - K = K^2$ (Расстояние до $x$ = расстояние до $y_1$, чтобы точки совпадали);

  \[
  1 - K = K^2 \Leftrightarrow K^2 + K - 1 = 0,  1 > K > 0 \Rightarrow K = \frac{\sqrt{5} - 1}{2} \approx 0.618
  \]

  Это равносильно этому условию
  \[
  \frac{l[a;x]}{l[a;y]} = \frac{l[a;y]}{l[a;b]}
  \]
  В этом равенстве как бы и говорится, что $[a;y]$ становится новым $[a;b]$, а $[a;x]$ должен стать новым $[a;y]$

  Собственно сам алгоритм:

  \begin{Verbatim}[fontsize=\small]
  K = (sqrt(5) -1 )/ 2
  x = b - K*(b-a); y = a + K*(b-a)
  fx = f(x); fy = f(y)

  while b - a > eps:
      if fx < fy: # минимум в [a;y]
          b = y
          y = x; fy = fx
          x = b - K*(b-a)
          fx = f(x)
      else:
          a = x
          x = y; fx = fy
          y = a + K*(b-a)
          fy = f(y)
  \end{Verbatim}

  Сходимость линейная, так как
  $| x_k - x^{*}| \le I_k$, где $I_k$~--- длина отрезка на $k$ итерации. А считается она так $I_k = K \cdot I_{k-1} = K^k I_0$ \\
  Итого $|x_k - x^{*}| \le I_k = K^k I_0$~--- \hyperref[note:DF2405BF-8B31-47EB-8DEF-B1A88034717C]{линейная скорость}
\end{method}
